\documentclass[11pt,a4paper]{article}

%====================================================
% Packages
%====================================================
\usepackage[margin=1in]{geometry}
\usepackage{amsmath,amssymb,amsfonts}
\usepackage{graphicx}
\usepackage{booktabs}
\usepackage{float}
\usepackage{natbib}
\usepackage{siunitx}
\usepackage{xcolor}
\usepackage{titlesec}
\usepackage{microtype}
\usepackage{setspace}
\usepackage{caption}

%====================================================
% Colours
%====================================================
\definecolor{sectionblue}{RGB}{31,78,121}
\definecolor{subsectionblue}{RGB}{46,95,138}
\definecolor{ruleblue}{RGB}{91,155,213}
\definecolor{citationblue}{RGB}{0,102,178}
\definecolor{urlblue}{RGB}{0,114,188}
\definecolor{textgray}{RGB}{80,80,80}

%====================================================
% Hyperlinks
%====================================================
\usepackage[
    colorlinks=true,
    linkcolor=sectionblue,
    citecolor=citationblue,
    urlcolor=urlblue
]{hyperref}

%====================================================
% Cross-References
%====================================================
\usepackage[
    nameinlink,
    capitalise,
    noabbrev
]{cleveref}

%====================================================
% Cross-Reference Names
%====================================================
\crefname{equation}{equation}{equations}
\Crefname{equation}{Equation}{Equations}

\crefname{figure}{figure}{figures}
\Crefname{figure}{Figure}{Figures}

\crefname{table}{table}{tables}
\Crefname{table}{Table}{Tables}

%====================================================
% Section Styling
%====================================================
\titleformat{\section}
  {\Large\bfseries\color{sectionblue}}
  {\thesection}
  {0.75em}
  {}
  [\color{ruleblue}\titlerule]

\titleformat{\subsection}
  {\large\bfseries\color{subsectionblue}}
  {\thesubsection}
  {0.75em}
  {}

\titleformat{\subsubsection}
  {\normalsize\bfseries\color{subsectionblue}}
  {\thesubsubsection}
  {0.75em}
  {}

\titlespacing*{\section}{0pt}{2.0em}{0.8em}
\titlespacing*{\subsection}{0pt}{1.5em}{0.5em}
\titlespacing*{\subsubsection}{0pt}{1.2em}{0.4em}

%====================================================
% General Formatting
%====================================================
\setstretch{1.08}

\setlength{\parindent}{0pt}
\setlength{\parskip}{0.55em}

\captionsetup{
    font=small,
    labelfont={bf,color=sectionblue}
}

%====================================================
% Custom Commands
%====================================================
\newcommand{\E}{\mathbb{E}}
\newcommand{\Q}{\mathbb{Q}}
\newcommand{\Var}{\operatorname{Var}}
\newcommand{\dd}{\,\mathrm{d}}

%====================================================
% Equation Numbering
%====================================================
\numberwithin{equation}{section}

%====================================================
% Title Information
%====================================================
\title{
    \vspace{-1.5cm}
    {\color{sectionblue}\rule{\textwidth}{1.2pt}}\\[0.6cm]
    \Huge\bfseries
    A Look at the Ho--Lee Short Rate Model\\[0.6cm]
    {\color{ruleblue}\rule{\textwidth}{0.6pt}}
}

\author{
    \large Godfrey Tshehla
}

\date{\today}

%====================================================
% Document
%====================================================
\begin{document}

\maketitle

\vspace{-0.5cm}

%====================================================
% Abstract
%====================================================
\begin{abstract}
\noindent
This paper studies the Ho--Lee short-rate model and its application to
term-structure modelling and interest-rate contingent-claim pricing.
The model is calibrated to an initial term structure through a
time-dependent drift, while interest-rate uncertainty is represented by
a constant volatility parameter. Zero-coupon bonds and European options
on zero-coupon bonds are valued analytically and numerically using Monte
Carlo simulation, finite-difference methods, and a calibrated lattice.
The numerical results are used to assess consistency across valuation
methods and to illustrate the implications of the model's Gaussian
dynamics and absence of mean reversion.
\end{abstract}


%====================================================
\section{Introduction}
%====================================================

Short-rate models provide a framework for describing the evolution of
interest rates and valuing securities whose cash flows depend on the term
structure. One of the earliest arbitrage-free term-structure models is
the Ho--Lee model introduced by \citet{HoLee1986}.

The Ho--Lee model specifies the instantaneous short rate under the
risk-neutral measure as a Gaussian process with a time-dependent drift
and constant volatility. Unlike equilibrium short-rate models, the drift
is selected so that the model is consistent with the observed initial
term structure.

The model is particularly useful as a benchmark because of its analytical
tractability and relatively simple calibration. Its principal limitation
is the absence of mean reversion, which causes the variance of the short
rate to increase without bound as the time horizon increases. The
Gaussian specification also permits negative interest rates.

The purpose of this paper is to examine the theoretical and numerical
properties of the Ho--Lee model. Particular attention is given to
term-structure calibration, short-rate dynamics, zero-coupon bond pricing,
and the valuation of a European option on a zero-coupon bond. Analytical
results are compared with Monte Carlo simulation, an implicit
finite-difference method, and a calibrated recombining lattice.


%====================================================
\section{The Ho--Lee Model}
%====================================================

Under the risk-neutral probability measure $\Q$, the instantaneous
short rate follows

\begin{align}
    \dd r_t
    &=
    \theta(t)\dd t
    +
    \sigma\dd W_t^{\Q}.
    \label{eq:ho_lee_sde}
\end{align}

Here, $\theta(t)$ is a deterministic time-dependent drift, $\sigma$ is
the constant short-rate volatility, and $W_t^{\Q}$ is a standard
Brownian motion under the risk-neutral measure.

Integrating \Cref{eq:ho_lee_sde} from $0$ to $t$ gives

\begin{align*}
    r_t
    &=
    r_0
    +
    \int_0^t \theta(s)\dd s
    +
    \sigma W_t^{\Q}.
\end{align*}

Since the Brownian component is Gaussian, the short rate is normally
distributed:

\begin{align*}
    r_t
    &\sim
    \mathcal{N}
    \left(
        r_0+\int_0^t\theta(s)\dd s,
        \sigma^2t
    \right).
\end{align*}

The first two moments are therefore

\begin{align*}
    \E^{\Q}[r_t]
    &=
    r_0+\int_0^t\theta(s)\dd s,
\end{align*}

and

\begin{align}
    \Var^{\Q}(r_t)
    &=
    \sigma^2t.
    \label{eq:ho_lee_variance}
\end{align}

As shown in \Cref{eq:ho_lee_variance}, the variance increases linearly
with the time horizon. This is a defining characteristic of the Ho--Lee
model and follows directly from the absence of a mean-reversion term.


%====================================================
\section{Term-Structure Calibration}
%====================================================

Let $P(0,T)$ denote the time-zero price of a zero-coupon bond with
maturity $T$. The instantaneous forward rate is defined as

\begin{align*}
    f(0,T)
    &=
    -\frac{\partial}{\partial T}
    \ln P(0,T).
\end{align*}

To ensure consistency with the initial term structure, the Ho--Lee drift
is specified as

\begin{align}
    \theta(t)
    &=
    \frac{\partial f(0,t)}{\partial t}
    +
    \sigma^2t.
    \label{eq:ho_lee_theta}
\end{align}

Consequently, $\theta(t)$ is not an independently estimated constant.
Instead, it is determined by the shape of the initial forward curve and
the assumed short-rate volatility.

For a continuously compounded zero yield $y(0,T)$, the initial discount
factor is

\begin{align*}
    P(0,T)
    &=
    \exp\left[-y(0,T)T\right].
\end{align*}

In the numerical implementation, a smooth parametric zero-yield curve is
used to construct the initial discount curve. The corresponding
instantaneous forward curve is obtained numerically, after which
$\theta(t)$ is evaluated using \Cref{eq:ho_lee_theta}.


%====================================================
\section{Zero-Coupon Bond Pricing}
%====================================================

The value at time $t$ of a zero-coupon bond paying one unit at maturity
$T$ is

\begin{align*}
    P(t,T)
    &=
    \E^{\Q}
    \left[
        \left.
        \exp
        \left(
            -\int_t^T r_s\dd s
        \right)
        \right|
        \mathcal{F}_t
    \right].
\end{align*}

Under the calibrated Ho--Lee model, the zero-coupon bond price is

\begin{align*}
    P(t,T)
    &=
    \frac{P(0,T)}{P(0,t)}
    \exp
    \left[
        -(T-t)
        \bigl(r_t-f(0,t)\bigr)
        -
        \frac{1}{2}
        \sigma^2t(T-t)^2
    \right].
\end{align*}

The bond price therefore admits the exponential-affine representation

\begin{align*}
    P(t,T)
    &=
    A(t,T)\exp[-B(t,T)r_t],
    \\
    B(t,T)
    &=
    T-t.
\end{align*}

At the initial time, the calibration condition implies

\begin{align*}
    P^{HL}(0,T)
    &=
    P^{market}(0,T).
\end{align*}

Thus, the model reproduces the initial term structure by construction.


%====================================================
\section{European Option on a Zero-Coupon Bond}
%====================================================

Consider a European option expiring at $T_1$ on a zero-coupon bond
maturing at $T_2$, where

\begin{align*}
    0
    &<
    T_1
    <
    T_2.
\end{align*}

For notional $N$ and strike price $X$, the call payoff at $T_1$ is

\begin{align*}
    C(T_1)
    &=
    N\max\left(P(T_1,T_2)-X,0\right).
\end{align*}

The corresponding put payoff is

\begin{align*}
    P_{\mathrm{put}}(T_1)
    &=
    N\max\left(X-P(T_1,T_2),0\right).
\end{align*}

For the Ho--Lee model, the integrated forward-bond volatility is

\begin{align*}
    \Sigma
    &=
    \sigma(T_2-T_1)\sqrt{T_1}.
\end{align*}

The time-zero call price is

\begin{align*}
    C(0)
    &=
    N
    \left[
        P(0,T_2)\Phi(d_1)
        -
        X P(0,T_1)\Phi(d_2)
    \right].
\end{align*}

The quantities $d_1$ and $d_2$ are given by

\begin{align*}
    d_1
    &=
    \frac{
        \ln
        \left(
            \dfrac{P(0,T_2)}
                  {X P(0,T_1)}
        \right)
        +
        \dfrac{1}{2}\Sigma^2
    }{\Sigma},
    \\
    d_2
    &=
    d_1-\Sigma.
\end{align*}

Here, $\Phi(\cdot)$ denotes the standard normal cumulative distribution
function.

The corresponding time-zero put price is

\begin{align*}
    P_{\mathrm{put}}(0)
    &=
    N
    \left[
        X P(0,T_1)\Phi(-d_2)
        -
        P(0,T_2)\Phi(-d_1)
    \right].
\end{align*}

The call and put prices satisfy put-call parity:

\begin{align*}
    C(0)-P_{\mathrm{put}}(0)
    &=
    N
    \left[
        P(0,T_2)-X P(0,T_1)
    \right].
\end{align*}

The at-the-money forward strike price is

\begin{align}
    X_{\mathrm{ATM}}
    &=
    \frac{P(0,T_2)}
         {P(0,T_1)}.
    \label{eq:atm_forward_strike}
\end{align}

The bond option is used primarily as a validation instrument for the
numerical implementation rather than as the central object of the
analysis.


%====================================================
\section{Numerical Implementation}
%====================================================

The analytical results are compared with three numerical approaches:
Monte Carlo simulation, an implicit finite-difference method, and a
calibrated recombining lattice.


%----------------------------------------------------
\subsection{Monte Carlo Simulation}
%----------------------------------------------------

The short-rate dynamics are discretised over a time interval $\Delta t$
according to

\begin{align*}
    r_{t+\Delta t}
    &=
    r_t
    +
    \theta(t)\Delta t
    +
    \sigma\sqrt{\Delta t}\,Z,
    \\
    Z
    &\sim
    \mathcal{N}(0,1).
\end{align*}

Given $M$ simulated short-rate paths, the zero-coupon bond price is
estimated by

\begin{align*}
    \widehat{P}(0,T)
    &=
    \frac{1}{M}
    \sum_{m=1}^{M}
    \exp
    \left(
        -\int_0^T r_s^{(m)}\dd s
    \right).
\end{align*}

For the European zero-coupon bond call option, the Monte Carlo estimator
is

\begin{align*}
    \widehat{C}(0)
    &=
    \frac{1}{M}
    \sum_{m=1}^{M}
    \exp
    \left(
        -\int_0^{T_1}r_s^{(m)}\dd s
    \right)
    N
    \max
    \left(
        P^{(m)}(T_1,T_2)-X,
        0
    \right).
\end{align*}


%----------------------------------------------------
\subsection{Finite-Difference Method}
%----------------------------------------------------

Let $V(t,r)$ denote the value of an interest-rate contingent claim. Under
the risk-neutral Ho--Lee dynamics, $V$ satisfies

\begin{align}
    \frac{\partial V}{\partial t}
    +
    \theta(t)\frac{\partial V}{\partial r}
    +
    \frac{1}{2}\sigma^2
    \frac{\partial^2 V}{\partial r^2}
    -
    rV
    &=
    0.
    \label{eq:ho_lee_pde}
\end{align}

The partial differential equation in \Cref{eq:ho_lee_pde} is solved
backward from the terminal payoff using an implicit finite-difference
scheme on a discretised short-rate grid.

For the zero-coupon bond, the terminal condition is

\begin{align*}
    V(T,r)
    &=
    N.
\end{align*}

For the European zero-coupon bond call option, the terminal condition at
the option expiry is

\begin{align*}
    V(T_1,r)
    &=
    N\max\left(P(T_1,T_2)-X,0\right).
\end{align*}


%----------------------------------------------------
\subsection{Calibrated Ho--Lee Lattice}
%----------------------------------------------------

A recombining short-rate lattice provides an alternative discrete
representation of the Ho--Lee dynamics. For a time increment $\Delta t$,
the short-rate spacing is

\begin{align*}
    h
    &=
    \sigma\sqrt{\Delta t}.
\end{align*}

At node $(n,j)$, the short rate is represented as

\begin{align*}
    r_{n,j}
    &=
    \alpha_n+(2j-n)h.
\end{align*}

The parameter $\alpha_n$ is calibrated at each time step so that the
state-price lattice reproduces the observed initial discount curve.

Let $Q_{n,j}$ denote the Arrow--Debreu state price associated with node
$(n,j)$. For an interior node, the state-price recursion can be written
as

\begin{align*}
    Q_{n+1,j}
    &=
    \frac{1}{2}
    Q_{n,j}
    \exp(-r_{n,j}\Delta t)
    +
    \frac{1}{2}
    Q_{n,j-1}
    \exp(-r_{n,j-1}\Delta t).
\end{align*}

The boundary nodes receive contributions from only one predecessor. The
calibration condition at time step $n+1$ is

\begin{align*}
    \sum_{j=0}^{n+1}Q_{n+1,j}
    &=
    P(0,t_{n+1}).
\end{align*}

For the European zero-coupon bond option, the lattice is constructed to
the bond maturity $T_2$. The underlying bond is first valued backward
from $T_2$ to the option expiry $T_1$. The call payoff at each expiry
node is

\begin{align*}
    C_{n_1,j}
    &=
    N\max\left(P_{n_1,j}(T_2)-X,0\right),
\end{align*}

where $n_1$ denotes the lattice time index corresponding to $T_1$. The
option values are subsequently discounted backward from $T_1$ to time
zero.


%====================================================
\section{Numerical Results}
%====================================================

The numerical experiments compare the analytical Ho--Lee solution with
Monte Carlo, finite-difference, and lattice valuations. The same initial
term structure and volatility parameter are used across all methods.


%----------------------------------------------------
\subsection{Initial Term Structure and Calibration}
%----------------------------------------------------

\Cref{fig:curve} presents the initial term structure used in the Ho--Lee
calibration.

\begin{figure}[H]
    \centering

    % Replace this placeholder after generating the figure:
    %
    % \includegraphics[
    %     width=0.78\textwidth
    % ]{figures/yield_curve.pdf}

    \fbox{
        \parbox[c][5cm][c]{0.78\textwidth}{
            \centering
            \color{textgray}
            Initial yield curve and calibration figure
        }
    }

    \caption{Initial term structure used to calibrate the Ho--Lee model.}
    \label{fig:curve}
\end{figure}


%----------------------------------------------------
\subsection{Short-Rate Dynamics}
%----------------------------------------------------

\Cref{fig:paths} illustrates simulated short-rate paths generated under
the calibrated risk-neutral dynamics.

\begin{figure}[H]
    \centering

    % Replace this placeholder after generating the figure:
    %
    % \includegraphics[
    %     width=0.78\textwidth
    % ]{figures/ho_lee_paths.pdf}

    \fbox{
        \parbox[c][5cm][c]{0.78\textwidth}{
            \centering
            \color{textgray}
            Simulated Ho--Lee short-rate paths
        }
    }

    \caption{Simulated short-rate paths under the Ho--Lee model.}
    \label{fig:paths}
\end{figure}


%----------------------------------------------------
\subsection{Zero-Coupon Bond Pricing}
%----------------------------------------------------

\Cref{tab:zcb_pricing} compares the zero-coupon bond valuation methods.

\begin{table}[H]
    \centering

    \caption{Comparison of Ho--Lee zero-coupon bond valuation methods.}
    \label{tab:zcb_pricing}

    \renewcommand{\arraystretch}{1.2}

    \begin{tabular}{lcc}
        \toprule
        \textbf{Method}
        &
        \textbf{Price}
        &
        \textbf{Pricing Error}
        \\
        \midrule
        Market
        & -- & -- \\
        Analytical
        & -- & -- \\
        Monte Carlo
        & -- & -- \\
        Finite Difference
        & -- & -- \\
        Lattice
        & -- & -- \\
        \bottomrule
    \end{tabular}
\end{table}


%----------------------------------------------------
\subsection{Zero-Coupon Bond Option Pricing}
%----------------------------------------------------

The option experiment uses an expiry $T_1$, bond maturity $T_2$, and
strike price $X$. When an at-the-money forward option is considered, the
strike price is determined from \Cref{eq:atm_forward_strike}.

\Cref{tab:option_pricing} compares the European zero-coupon bond option
valuation methods.

\begin{table}[H]
    \centering

    \caption{Comparison of Ho--Lee zero-coupon bond option valuations.}
    \label{tab:option_pricing}

    \renewcommand{\arraystretch}{1.2}

    \begin{tabular}{lcc}
        \toprule
        \textbf{Method}
        &
        \textbf{Price}
        &
        \textbf{Error vs.\ Closed Form}
        \\
        \midrule
        Closed Form
        & -- & -- \\
        Monte Carlo
        & -- & -- \\
        Finite Difference
        & -- & -- \\
        Lattice
        & -- & -- \\
        \bottomrule
    \end{tabular}
\end{table}

The analytical values provide natural benchmarks against which the
numerical methods can be assessed. Monte Carlo errors arise from sampling
uncertainty and time discretisation, whereas finite-difference and
lattice errors primarily reflect discretisation of the state and time
dimensions.


%====================================================
\section{Volatility Sensitivity and Model Behaviour}
%====================================================

The volatility parameter $\sigma$ is the principal free model parameter
once the initial term structure has been specified. It directly controls
the dispersion of future short rates:

\begin{align*}
    \Var^{\Q}(r_t)
    &=
    \sigma^2t.
\end{align*}

Increasing $\sigma$ therefore widens the distribution of future short
rates. The effect is particularly visible in the value of the
zero-coupon bond option because its integrated volatility is

\begin{align*}
    \Sigma
    &=
    \sigma(T_2-T_1)\sqrt{T_1}.
\end{align*}

The sensitivity of an option value $V$ to the Ho--Lee volatility can be
summarised by its Vega:

\begin{align*}
    \mathrm{Vega}
    &=
    \frac{\partial V}{\partial\sigma}.
\end{align*}

A numerical approximation can be obtained through a symmetric volatility
bump:

\begin{align*}
    \mathrm{Vega}
    &\approx
    \frac{
        V(\sigma+\Delta\sigma)
        -
        V(\sigma-\Delta\sigma)
    }{
        2\Delta\sigma
    }.
\end{align*}


%====================================================
\section{Limitations of the Ho--Lee Model}
%====================================================

Despite its tractability, the Ho--Lee model has several important
limitations.

First, there is no mean-reversion mechanism in the short-rate dynamics.
As established in \Cref{eq:ho_lee_variance},

\begin{align*}
    \Var^{\Q}(r_t)
    &=
    \sigma^2t.
\end{align*}

Consequently,

\begin{align*}
    \lim_{t\rightarrow\infty}
    \Var^{\Q}(r_t)
    &=
    \infty.
\end{align*}

Short-rate uncertainty therefore increases indefinitely with the time
horizon.

Second, the Gaussian specification implies

\begin{align*}
    \Pr(r_t<0)
    &>
    0.
\end{align*}

Thus, for suitable parameter values and horizons, negative short rates
occur with positive probability.

Finally, the model assumes a constant instantaneous volatility $\sigma$.
Although this provides substantial analytical convenience, it restricts
the model's ability to reproduce more complex volatility structures
observed in interest-rate markets.

These limitations motivate extensions of the Gaussian short-rate
framework that introduce mean reversion while preserving consistency
with the initial term structure \citep{BrigoMercurio2006}.


%====================================================
\section{Conclusion}
%====================================================

The Ho--Lee model provides a simple and tractable arbitrage-free framework
for modelling the term structure of interest rates. Its deterministic
time-dependent drift allows the model to reproduce an observed initial
discount curve, while its Gaussian dynamics lead to analytically
tractable distributions and bond prices.

The numerical implementation considered analytical valuation, Monte
Carlo simulation, an implicit finite-difference method, and a calibrated
recombining lattice. Comparing these methods provides a useful numerical
validation of the model and illustrates the different sources of
discretisation and simulation error.

The model's simplicity is also the source of its principal limitations.
In particular, the absence of mean reversion implies that short-rate
variance increases linearly with time, while the Gaussian distribution
permits negative rates. Nevertheless, the Ho--Lee model remains an
important benchmark for understanding arbitrage-free term-structure
modelling and provides a natural foundation for more sophisticated
short-rate models.


%====================================================
% References
%====================================================

\bibliographystyle{plainnat}
\bibliography{references}

\end{document}